\documentclass[12pt]{article}

\usepackage[a4paper, total={6in, 8in}]{geometry}
\usepackage{textcomp}

\begin{document}
\setlength{\parindent}{0pt}

\def \t {\textrightarrow}
\def \v {\vspace{1em}}
\def \bi {\begin{itemize}}
\def \ei {\end{itemize}}
\def \s[#1] {\section*{#1}}
\def \ss[#1] {\subsection*{#1}}
\def \sss[#1] {\subsubsection*{#1}}

\s[Svevo]
La psicanalsi viene dopo la letteratura \t l'indagine dell'io parte dalla letteratura \\
Due pilastri che fanno la svolta del decadentismo \\
Svevo vive una realtà di frontiera, a Trieste \t città del dualismo \t e continua provocazione culturale che è sollecitata dal suo dualismo interiore \\
Lui vive questa realtà nel nome: Ettore Schmidt e Italo Svevo \t duplice natura \\
Filone in cui tutte le provocazioni esterne vengono riassorbite nell'animo dell'artista \t è la precedenza di una nuova visione interna del romanzo, che è agli antipodi rispetto al romanzo manzoniano e anche verghiano \\
Se rimane qualche eco del mondo verista, questa è relegata alla descrizione degli ambienti nel romanzo "Selinità" \\
Il focus è sul personaggio e la sua anima \t narrazione soggettiva, e uno smembramento della struttura del romanzo \\
Equilibrio iniziale e poi rottura di questo equilibrio \t in Pirandello anche il romanzo rimane aperto, e lascia la possibilità di siglare una conclusione \t e anche il teatro \\
Il romanzo è completamente frammentato, proprio come l'io è diviso \t Virginia Woolf muore nel 41 \t il romanzo moderno \\
Il romanzo moderno ha prospettive nuove, di una voce che non rinuncia solo agli intrecci mozzafiato, ma ha nel suo comporre a rappresentare un universo conflittuale senza preclusioni, e a dare voce al proprio spirito mutevole \\
L'essenza dell'opera moderna sono le oscure zone dell'io \t mostra il mondo oscuro dell'anima umana \\
"Ulixis" è presente \t la mente riceve una miriade di impressioni \t come fa ad avere linearità? \\
Mette in evidenza la fragilità interiore \\
Romanzo di pirandello è romanzo saggio \\
La malattia fisica è occasione per il dolore interiore \t letteratura tedesca influenza molto

\v

Trionfo della donna, le ampie riflessioni filosofiche che caratterizzano la voce narrante della coscienza di zeno \t la montagna incantata di Thomas Mann è un punto di riferimento \\
Anche la rovinosa caduta dell'"Uomo senza qualita" \t parabola involutiva dell'essere umano \t inabilità a vivere e quindi inettitudine \\
Marcel Proust con la recerce \t romanzo senza trama è quello di Joyce \t quando si parla di paralisis parlo di Joyce \\
Quando parlo di libere associazioni, parlo anche di Italo svevo \t romanzo si costruisce sulla terapia e sulle associazioni libere \\
Woolf prende i fatti, e dai fatti l'interiorizzazione \t Zeno Cosini parte invece dalla memoria che fa riaffiorare, e costruisce la sua interiorizzazione \\
Romanzo della crisi, crisi della persona

\ss[Vita]
Poliglottismo di Trieste, perchè ci sono tanti linguaggi che si parlano \\
Un triestino eredita il rigore asburgico, ma anche la leggerezza italian \t è una città di frontiera, non si trova mai in stanzialità fissa \\
Duplicità ha ricaduta su Svevo \t testimonianza di Gianni Stufaric, che ricorda Svevo come animatore brillante delle compagnie di intellettuali presso il bar Garibaldi \\
I caffe erano i luoghi del ritrovo \t Svevo è socievole, di mondo, affabile, ma allo stesso tempo tormentato e angosciato \\
Ha tempra di scrittore europeo \t il nodo commerciale che Trieste rappresenta come imperso Asburgico \t nel tempo diventa anche di industrie, che è quella del profitto \t e Svevo sente passione per la letteratura \\
Vita rutinaria senza scossoni \t si impiegherà come dipendente nella fabbrica del soucero \t si era sposato con Lidia Veneziani \\
Lei è la donna che vede che lui non ha un occhio storto? \\
L'industria è di vernici \t questo lo porta a mettere in pratica i suoi studi, che non erano stati classici \t frequenta istituto tecnico chimico, quindi non ha formazione classica \\
Avvilente condizione, lui vuole fare altro \t si adatta, e schiacchia ma non sopprime la sua passione per la scrittura \\
Si ricoda la nuova disciplina della psicanalisi, che entra nel dibattito grazie ai letterati \t questa ha presa nell'ambiente triestino, dove è presente il migliore allievo di Freud: Eduard Weiss \\
Questa entra nella vita di Svevo, il cognato aveva una distorsione della personalita \t quindi lui conosce sul campo la pratica \\
Giorgio Voghera, scrittore triestino \t ?? \\
Zeno cosini dimostra che la psicanalisi fallisce ??

\v

Chiedere questa parte \\
La ipnosi \t tipo ti terapia che voleva curare l'isteria \\
Si parla di una parascienza \t la sua esperienza di letteratura anticipa \t svevo non salva la coscienza \t cos'è allora il mezzo per approdare alla coscienza? è la scrittura, che è conoscenza di se \\
Eugenio Montale, nel suo omaggio a Italo Svevo \t parla di un desiderio di sondare nel profondo dei meandri dell'anima, che sono al di la del fenomeno, ma sono in uno strato ancora + sotterraneo \\
Montale legge Svevo \t e ne rimane colpito \\
La letteratura come vita, di Bo \t letteratura come vita, la società moderna con il suo drammatico correre \t si trova quindi un analisi chiarificante della vita dell'uomo \\
E critica ai filosofi \t che avrebbero portato a galla l'animo umano \\
Alfonso Niti nel romanzo ? \t Darwin, la societa mangia lui che si suicida \t selezione naturale \\
La critica di Marx e Spencer, Svevo trova conflittualità nel pensiero di Marx \t che indubbiamente ha un riflesso su Svevo quando pensa alla condizione psicologica dei personaggi, succubi dell'industrializzazione moderna \t che lui vive \\
Lui si trasferira anche a Roma \\
Chiedere critica a Marx che fa Svevo?? \\
Critica a Schopenahuer: da dove viene questa sofferenza? lui non si sofferma sull'orgine ma sulla situazione di dolore \\
Come se fosse un apparato che ha già deciso per noi
\end{document}
